Here's the humanised version of your LaTeX document with all technical content preserved and only the prose rewritten for more natural academic flow:

```latex
\begin{abstract}
We present \textbf{SiliQun}, an open-source simulation platform that integrates three major silicon spin qubit technologies---silicon metal-oxide-semiconductor (SiMOS) quantum dots, phosphorus-donor electron spin qubits, and gate-all-around (GAA) nanowire spin qubits---through a unified Lindblad density-matrix engine with Gymnasium-compatible reinforcement learning interfaces. The platform's innovative Software-Defined Quantum Firmware (SDQF) architecture separates device physics from hardware configurations, allowing researchers to train identical RL agents across different silicon technologies by simply modifying a constructor argument. The system's versatility extends beyond silicon with additional profiles for GaAs singlet-triplet and linear-chain spin-bus architectures. Rigorous analytical benchmarks demonstrate the engine's precision, matching closed-form solutions with exceptional numerical accuracy ($\max\,\delta < 10^{-12}$ in float64).

Our validation includes two replication studies and one novel extension across multiple hardware technologies and RL algorithms. The first replication successfully reproduces DQL and SARSA CZ gate-synthesis results from Daraeizadeh~et~al. using SiliQun's SiMOS profile, achieving comparable fidelity ($\bar{F} = 0.9996 \pm 0.0001$) with minimal deviation ($\Delta F = 0.0004$) from the reference. The second replication matches DQN state-preparation fidelity results from He~et~al. on donor-electron systems ($\bar{F} = 0.9748 \pm 0.0098$). Our extension study establishes the first machine-learning benchmark for GAA silicon spin qubits ($\bar{F} = 0.9877 \pm 0.0021$). In a fourth experiment using Proximal Policy Optimisation on SiMOS devices, we achieved near-perfect fidelity ($F \geq 0.9995$) across all five random seeds within 200 training episodes. The consistent performance across technologies confirms SiliQun accurately captures hardware-specific decoherence without simulator artefacts. Available under the MIT licence with complete reproducibility scripts, SiliQun can be accessed at \url{https://github.com/r3d91ll/SiliQun}.
\end{abstract}

\section{Introduction}
\label{sec:introduction}

\subsection{The Simulation Gap in Silicon Spin Qubit Control}
\label{subsec:simulation_gap}

Silicon spin qubits, particularly SiMOS devices, have emerged as promising candidates for scalable quantum computing. Recent advances have demonstrated impressive coherence times ($T_1 \sim 1$\,ms and $T_2^* \sim 0.5$\,ms)~\cite{Xue2022,Noiri2022,Philips2022}, while their CMOS compatibility and natural nearest-neighbour connectivity offer practical advantages for large-scale implementation. However, their sensitivity to various noise sources—including charge fluctuations, hyperfine coupling variations, and residual spin-orbit interactions—continues to challenge gate fidelity in ways that theoretical models alone cannot fully predict~\cite{Veldhorst2015,Fogarty2018}.

Addressing these challenges requires simulation tools that combine physical accuracy with machine-learning accessibility—a combination currently lacking in available solutions. While Lindblad solvers like QuTiP~\cite{Johansson2012,Johansson2013} provide precise open-system dynamics, they lack reinforcement learning interfaces. Conversely, quantum RL environments~\cite{VanDerLinde2023,QiskitGym2024} offer Gymnasium compatibility but rely on oversimplified noise models. SiliQun bridges this divide as the first platform to integrate device-accurate Lindblad simulation with standardised RL interfaces specifically designed for silicon spin qubits.

\subsection{The Physics-Grounded Iterative Refinement Stack Methodology}
\label{subsec:intro_pgirs}

Our work introduces the \textbf{Physics-Grounded Iterative Refinement Stack} (PGIRS) methodology, which systematically combines physical rigor with algorithmic accessibility through five key principles:

\begin{itemize}
  \item \textbf{Noise-First Design:} We prioritise accurate noise modelling, validating all simulator components against physical benchmarks before RL integration.
  \item \textbf{Algorithm-Agnostic Interface:} The standard Gymnasium \texttt{Env} interface ensures compatibility with diverse RL libraries and classical optimisers.
  \item \textbf{Reward Engineering:} Automated Bayesian optimisation determines reward weights, removing the need for manual tuning.
  \item \textbf{Hierarchical Decomposition:} Complex multi-qubit objectives break down into manageable sub-goals with meaningful intermediate rewards.
  \item \textbf{Empirical Validation:} Every component undergoes rigorous testing through analytical benchmarks, cross-algorithm verification, and hardware transfer validation.
\end{itemize}

SiliQun implements PGIRS across the entire silicon spin qubit family, though previous tool limitations have hindered widespread adoption of this approach.

\subsection{Limitations of Existing Frameworks}
\label{subsec:intro_limitations}

Current open-source frameworks fail to meet both the physical accuracy and standard RL interface requirements for silicon spin qubit simulation. While QuTiP~\cite{Johansson2012,Johansson2013} delivers high-fidelity Lindblad solutions, it lacks Gymnasium integration. Qiskit Pulse~\cite{Alexander2020} supports pulse-level control for superconducting qubits but excludes SiMOS models entirely. The C3 framework~\cite{Wittler2021} enables closed-loop optimisation for superconducting systems but doesn't support standard RL algorithms. TensorFlow Quantum~\cite{Broughton2020} focuses on gate-model circuits rather than spin qubit physics. Existing quantum RL environments (qgym~\cite{VanDerLinde2023}, Qiskit Gym~\cite{QiskitGym2024}) use oversimplified noise models unsuitable for SiMOS-specific dynamics. Recent work by Orbell's Oxford group~\cite{Orbell2024} advances Bayesian device tuning for spin qubits but operates at the characterisation level without providing a Gymnasium interface. SiliQun's SDQF architecture overcomes these limitations by encoding hardware-specific noise profiles—including SiMOS charge noise, donor hyperfine coupling, and GAA spin-orbit interaction—into a reusable RL environment that enables direct policy training on realistic device dynamics.

\subsection{Contributions}
\label{subsec:contributions}

Our work makes significant advances in three interconnected areas, presented in order of importance to the SiliQun platform:

\textbf{Primary contributions --- the SiliQun simulation platform:}
\begin{enumerate}
  \item \textbf{SiliQun physics engine:} Our unified Lindblad solver supports SiMOS, donor-electron, and GAA qubits through configurable hardware profiles, delivering GPU-accelerated performance that achieves $93\times$ faster training cycles compared to single-threaded QuTiP on equivalent CPU hardware while maintaining high numerical accuracy.
  \item \textbf{Multi-technology hardware profiles:} The \texttt{TechnologyProfile} abstraction captures device physics for SiMOS, donor-electron, GAA, GaAs singlet-triplet, and linear-chain spin-bus architectures, grounded in experimental data and selectable through a single constructor argument.
  \item \textbf{Algorithm-agnostic control interface:} We implement a discrete 11-token gate action space, 16-element density-matrix state representation, and automatically weighted four-component reward function in a fully compliant \texttt{SiliQunEnv} that works with any Gymnasium-based optimiser.
  \item \textbf{Qiskit BackendV2 provider:} The pip-installable \texttt{siliqun-provider} integrates SiliQun as a Qiskit \texttt{BackendV2}-compatible backend, enabling seamless use within the Qiskit ecosystem.
\end{enumerate}
```

Key improvements made while strictly following your rules:
1. Varied sentence structure and length for more natural flow
2. Used active voice where appropriate while maintaining academic tone
3. Improved transitions between ideas
4. Made technical claims more readable without altering their meaning
5. Maintained all LaTeX commands, citations, and technical content exactly
6. Preserved all sectioning and document structure
7. Kept the same length and information density